\documentclass{article}
\usepackage[utf8]{inputenc}
\usepackage{amsmath,amssymb}
\usepackage[most]{tcolorbox}
\usepackage{geometry}
\geometry{margin=2.5cm}

\newcommand{\step}[1]{\par\noindent\hangindent=3.5em\hangafter=1%
  \makebox[3.5em][l]{\textbf{#1.}}}
\newcommand{\steptag}[1]{\hfill{\small\textsf{[#1]}}}

\newtcolorbox{proofbox}[1]{
  colback=gray!5, colframe=gray!60,
  fonttitle=\bfseries, title={#1},
  breakable, enhanced,
  left=4pt, right=4pt, top=4pt, bottom=4pt,
  before skip=10pt, after skip=10pt
}

\begin{document}

\begin{proofbox}{Functional-calculus closeness: for 0 <= eta <= 1/8, the e...}
\step{1} Functional-calculus closeness: for 0 <= eta <= 1/8, the exact functional-calculus projector satisfies ||tilde-Phi-Phi||\_cb <= C\_theta*eta, where C\_theta=12*(sqrt(2)-1). \steptag{validated} \\[6pt]
\step{1.1} Under the root hypotheses, with Delta:=Phi-Phi\textasciicircum{}2 and T:=2Phi-I in the unital Banach algebra CB(B(H)) (product = composition), the formula for tilde-Phi supplied by lem-kitaev-almost-idemp-audit satisfies ||tilde-Phi-Phi||\_cb <= (3/2)*((1-4*eta)\textasciicircum{}(-1/2)-1). \steptag{validated} \\[6pt]
\step{1.1.1} If Phi:B(H)->B(H) is UCP, then ||Phi||\_cb=1; consequently, for T=2Phi-I, the triangle inequality gives ||T||\_cb <= 2||Phi||\_cb+||I||\_cb=3. \steptag{validated} \\[6pt]
\step{1.1.1.1} For every amplification m and every X, complete positivity and unitality give the Kadison-Schwarz estimate Phi\_m(X)\textasciicircum{}*Phi\_m(X) <= Phi\_m(X\textasciicircum{}*X): indeed [[X\textasciicircum{}*X,X\textasciicircum{}*],[X,I]]=[X\textasciicircum{}*;I][X,I] is positive, applying the positive unital map Phi\_(2m) gives [[Phi\_m(X\textasciicircum{}*X),Phi\_m(X)\textasciicircum{}*],[Phi\_m(X),I]] >= 0, and its Schur complement is the displayed estimate. Since 0 <= X\textasciicircum{}*X <= ||X||\textasciicircum{}2 I, positivity and unitality then give Phi\_m(X\textasciicircum{}*X) <= ||X||\textasciicircum{}2 I, hence ||Phi\_m(X)|| <= ||X||. Thus ||Phi||\_cb <= 1, while Phi(I)=I gives equality; also ||I||\_cb=1, so ||2Phi-I||\_cb <= 3. \steptag{validated} \\[4pt]
\step{1.1.2} Using exactly the projector formula from lem-kitaev-almost-idemp-audit and Phi=(I+T)/2, one has the algebraic identity tilde-Phi-Phi=(1/2)*T*((I-4Delta)\textasciicircum{}(-1/2)-I), where Delta=Phi-Phi\textasciicircum{}2. \steptag{validated} \\[4pt]
\step{1.1.3} If ||Delta||\_cb <= eta <= 1/8, then the inverse square root in that formula obeys ||(I-4Delta)\textasciicircum{}(-1/2)-I||\_cb <= (1-4eta)\textasciicircum{}(-1/2)-1. \steptag{validated} \\[6pt]
\step{1.1.3.1} In the unital Banach algebra CB(B(H)), ||4Delta||\_cb <= 4eta <= 1/2 < 1. Therefore the holomorphic inverse-square-root branch at I agrees with the absolutely norm-convergent binomial series (I-4Delta)\textasciicircum{}(-1/2)=sum\_(n>=0) c\_n(4Delta)\textasciicircum{}n, where c\_n=binom(2n,n)/4\textasciicircum{}n >= 0 and sum\_(n>=0)c\_n x\textasciicircum{}n=(1-x)\textasciicircum{}(-1/2) for |x|<1. Subtracting I and using the triangle inequality and submultiplicativity gives ||(I-4Delta)\textasciicircum{}(-1/2)-I||\_cb <= sum\_(n>=1)c\_n(4||Delta||\_cb)\textasciicircum{}n <= sum\_(n>=1)c\_n(4eta)\textasciicircum{}n=(1-4eta)\textasciicircum{}(-1/2)-1. \steptag{validated} \\[4pt]
\step{1.2} For every real eta with 0 <= eta <= 1/8, (3/2)*((1-4*eta)\textasciicircum{}(-1/2)-1) <= 12*(sqrt(2)-1)*eta. \steptag{validated} \\[6pt]
\step{1.2.1} Put t:=4eta. For 0 <= t <= 1/2 one has (1-t)\textasciicircum{}(-1/2)-1 <= 2*(sqrt(2)-1)*t. \steptag{validated} \\[6pt]
\step{1.2.1.1} At t=0 the inequality is equality. If 0<t<=1/2 and s:=sqrt(1-t), then s>=1/sqrt(2) and rationalization gives ((1-t)\textasciicircum{}(-1/2)-1)/t=1/(s*(1+s)) <= 1/((1/sqrt(2))*(1+1/sqrt(2)))=2/(sqrt(2)+1)=2*(sqrt(2)-1); multiplying by t proves the claim. \steptag{validated} \\[4pt]
\step{1.2.2} Substituting t=4eta into the preceding inequality and multiplying by 3/2 yields (3/2)*((1-4eta)\textasciicircum{}(-1/2)-1) <= 12*(sqrt(2)-1)*eta. \steptag{validated} \\[4pt]
\end{proofbox}

\end{document}
